% =====================================================================
%  Partial topology identification and cost-aware probing for
%  low-voltage feeders: method and algorithms
%  Paper-style draft v17, 2026-11-27
%
%  Status of the parts (read before quoting anything from this file):
%   * Both algorithms are implemented in topology_study/ptid/method
%     (self-contained: numpy, scipy and gurobipy; development only).
%     Development results: topology_study/RESULTS.md. Algorithm 2
%     (probe.py) ran on simulated days only, with a placeholder psi: NOT
%     evaluated.
%   * The registered sealed test evaluates the predecessor of Algorithm 1
%     (identification/final.py at commit 71684ff); Section VII lists the
%     differences. Its result does not carry over to this method.
%  History: v4 f5299d4, v5 5e36b4c, v6 4881a92, v7 dd8cd4a, v8 52313c3 and
%  b1c5e4d, v9 7bbe494 (load gate, meter-level labels), v10 91e80a0
%  (proofs repaired after a review). v11 (user, 2026-11-26): only results
%  the method needs, each with an elementary proof. Removed: the glue and
%  removal-gap lemmas, the asymptotic separation theorem and its corollary
%  (no consistency claim is made; reliability is measured on data), the
%  evidence proposition (now two lines of text), the parts of the reach
%  lemma that needed extra conditions (now an example in the text).
%  v12 (user, 2026-11-26): the guarantee the user proposed is back, in an
%  elementary, non-asymptotic form: Theorem 1 (removing a true line costs
%  the true signal the remaining tree cannot reproduce, minus noise in a few
%  directions; a redundant line gains only noise) and Corollary 1 (when
%  these hold, the true tree is the best, its groupings are decided, and the
%  anchored search keeps it). Proof: two facts about least squares and
%  Pythagoras for projections. The other results are labelled bookkeeping.
%  v13 (2026-11-26): aligned with the code (identification_v8): the flip
%  test of a grouping outside the tree (a support of it required), the
%  evidence interval's upper end from a re-solve over the pool, the
%  certificate scope in the output, the planner choices of probing and a
%  last gate after hour 23, and the wording of rho and of Appendix B.
%  v14 (user questions, 2026-11-26): the load-profile projection D1 is a
%  set-up step only (sigma^2 and Appendix B); eq. (5) is now the saturated
%  model. Theorem 1(b) bounds the gain of extra lines by the noise along the
%  directions they add (the old bound used all of the tree's directions, so
%  Corollary 1's condition (B) could almost never hold), with the BIC
%  reason stated. Pricing enumerates nodes up to K children and uses the
%  node MIQP beyond (was 12), and Algorithm 1 names it.
%  v15 (solver fix, 2026-11-26): on 47-60 meter feeders the tree MIQP at
%  forced 1e-9 tolerances declared feasible relaxations infeasible and
%  claimed false optima, and the local search stopped at its time limit far
%  from a local optimum. Now: default tolerances; local search to
%  convergence (every candidate line's value checked) before and after the
%  MIQP; every solver claim checked against the trees known; the MIQP's
%  coefficients scaled to the size of their columns (unit-diagonal Gram).
%  v16 (2026-11-27, the review of the whole project): the code moved to
%  topology_study/ptid; the node MIQP's coefficients are scaled like the
%  tree MIQP's; probing checks an outside target's flip test against its
%  best local swap; the status paragraph follows the development runs.
%
%  Compile with pdfLaTeX (e.g. on Overleaf). No external .bib needed.
% =====================================================================
\documentclass[journal,onecolumn,11pt]{IEEEtran}

\usepackage{amsmath,amssymb}
\usepackage{booktabs}
\usepackage{array}
\usepackage{algorithm}
\usepackage[noend]{algpseudocode}
\usepackage{url}
\usepackage[hidelinks]{hyperref}

\newtheorem{theorem}{Theorem}
\newtheorem{corollary}{Corollary}
\newtheorem{lemma}{Lemma}
\newtheorem{proposition}{Proposition}
\newtheorem{definition}{Definition}
\newtheorem{remark}{Remark}

\newcommand{\tp}{^{\top}}
\newcommand{\one}{\mathbf{1}}
\newcommand{\Lin}{\mathcal{L}}
\newcommand{\Nset}{\mathcal{N}}
\newcommand{\Bset}{\mathcal{B}}
\newcommand{\Sset}{\mathcal{S}}
\newcommand{\Tset}{\mathcal{T}}
\newcommand{\Uset}{\mathcal{U}}
\newcommand{\Aset}{\mathcal{A}}
\newcommand{\Kset}{\mathcal{K}}
\newcommand{\Gset}{\mathcal{G}}
\newcommand{\cross}{\pitchfork}
\algrenewcommand\algorithmicrequire{\textbf{Input:}}
\algrenewcommand\algorithmicensure{\textbf{Output:}}

\begin{document}

\title{Partial Topology Identification and Cost-Aware Probing\\ for Low-Voltage Feeders}
\author{Method draft v17 of 27~November~2026 (authors to be added)}
\maketitle

\begin{abstract}
Smart-meter voltage and power records can support a group of downstream customers even when several feeder trees fit the
data. We fit bounded linearised power-flow models over laminar families of downstream sets. When the measured power
balance indicates an unmetered load, its estimated profile enters the fit as an additional leaf whose location varies
across candidate trees. For each grouping in a reference tree, we ask whether a rival without that grouping can match the
reference within a specified score margin. A feasible rival demonstrates ambiguity relative to the reference. A valid
exclusion certifies agreement among all near-optimal trees in the specified pool. Unresolved searches retain a separate
label. Candidate generation uses starting trees, conditional enumeration under stable groupings, and single-line pricing.
An exhaustive check on four small inputs examines sensitivity to the pool. A separate probing procedure ranks feasible
demand-response actions by their forecast ability to distinguish rivals and their cost.
\end{abstract}

\begin{IEEEkeywords}
Topology identification, smart meters, partial identification, laminar families, mixed-integer quadratic programming,
unmetered load, active probing, demand response.
\end{IEEEkeywords}

\noindent\textbf{Evaluation status.} The four-input audit in Section~\ref{sec:tests} checks the present identification
implementation. An earlier version returned the registered predecessor's tree in 234 of 240 development cases with at
most 25 meters. The larger historical runs and the predecessor's sealed test are described separately in the project
results. Probing has been exercised on simulated days; its cost--coverage benefit remains to be evaluated.

% ---------------------------------------------------------------------
\section*{Nomenclature}
\begin{tabular}{@{}p{0.27\textwidth}p{0.68\textwidth}@{}}
\toprule
$\Nset=\{1,\dots,n\}$, $t$, $N$ & Metered customers, time slots, rows of data ($N=Tn$)\\
$s$; $\Nset^+$ & Unmetered load, a leaf with estimated power and no voltage; all leaves ($\Nset\cup\{s\}$ or $\Nset$)\\
$\varsigma$, $\varsigma_0$ & Unmetered share of the metered power (\ref{eq:gate}); gate threshold (1\%)\\
$y_{i,t}=V_{0,t}^2-V_{i,t}^2$ & Squared-voltage drop at meter $i$\\
$C\subseteq\Nset^+$; $z_C$, $h_C$; $C\cross D$ & Support of a line; its meters, whether it carries the load; $C$ and $D$ cross\\
$\pi(C)=C\cap\Nset$; $\Gset(F)$ & Meters of a support; the groupings of a family\\
$F\supseteq\Bset$, $\Lin(F)$, $k(F)$ & Topology as a laminar family; its lines; their number\\
$\Omega_0$; $\Omega$ & Initial candidate lines (Appendix~\ref{app:pool}); current pool\\
$H(F)$; $H^{\rm f}_M$ & Score (\ref{eq:score}), half-BIC units; flip problem (\ref{eq:flip})\\
$m=4.6$ & Score margin of labels and anchors\\
$e_M$, $F_M$ & Gap of grouping $M$ (\ref{eq:gap}) and its rival tree\\
$\Aset$; $\Kset_A$, $k_A$; $K$ & Anchors; units of anchor $A$ and their number; largest node searched by unions\\
$\Tset$; $u$, $\Uset_h$, $c(u)$; $\psi$ & Probing target set; block, feasible blocks, price; value of one unit of evidence\\
\bottomrule
\end{tabular}

% ---------------------------------------------------------------------
\section{Model}\label{sec:model}

A feeder has a root, the point of common coupling (PCC), and $n$ metered customers. Meter $i$ reports $p_{i,t}$,
$q_{i,t}$, $V_{i,t}$; the PCC reports $V_{0,t}$ and, where metered, $P^{\rm pcc}_t$, $Q^{\rm pcc}_t$. Every line feeds a set
of leaves, its \emph{support}; the supports of a tree are laminar (any two are nested or disjoint), and a laminar family
containing $\Bset=\{\{1\},\dots,\{n\},\Nset^+\}$ fixes the tree up to internal nodes without leaves, which voltages cannot
see. A topology is a laminar family $F\supseteq\Bset$ over $\Nset^+$ with lines $\Lin(F)=F\setminus\Bset$. With the
linearised DistFlow equations,
\begin{equation}\label{eq:model}
y_t=\sum_{C\in F} z_C\big[r_C\,(z_C\tp p_t+h_C u^P_t)+x_C\,(z_C\tp q_t+h_C u^Q_t)\big]+\varepsilon_{0,t}\one+o+e_t ,
\end{equation}
with $z_C\in\{0,1\}^n$ the meters of $C$, $h_C=1$ if $s\in C$, $0\le r_C,x_C\le\bar\theta$ twice the line's resistance and
reactance in pu ($\bar\theta=2$ never binds), $\varepsilon_{0,t}$ the PCC voltage error, $o$ fixed meter offsets and $e_t$
meter noise. One unmetered load is allowed; its power is the PCC balance minus the losses estimated from the meters' own
voltages,
\begin{equation}\label{eq:balance}
u^P_t=P^{\rm pcc}_t-\sum_i p_{i,t}\Big(1+\frac{V_{0,t}-V_{i,t}}{V_{i,t}}\Big),\qquad u^Q_t\ \text{likewise}.
\end{equation}
When the gate of Section~\ref{sec:gate} finds a load, it is one more leaf $s$, with its estimated balance profile held fixed and no voltage reading;
otherwise there is no leaf $s$. The supports that carry the load all contain $s$, so in a laminar family they are nested:
they are the load's path (Proposition~\ref{prop:path}). The load can sit below any line, below the trunk, or at the PCC
(beside the trunk; the meters' trunk is then the line $\Nset$ without $s$). Its own branch $\{s\}$ affects no voltage: its
columns are zero and it carries no price. A line is supported by how the loads on one side of its branching point move the
voltages on the other side; a meter's own cable has no such branching evidence, so only lines with at least two leaves are
topology decisions. The data come from AC power flow; (\ref{eq:model}) is used inside the method only. \emph{Assumptions
stated, not modelled:} the loss estimate is approximate, several unmetered loads are not one aggregate, and errors in the
measured $p,q$ enter the regressors.

% ---------------------------------------------------------------------
\section{Is there an unmetered load?}\label{sec:gate}

Before any topology is fitted, the PCC balance decides whether a load is modelled at all:
\begin{equation}\label{eq:gate}
\varsigma=\frac{\sum_tu^P_t}{\sum_t\sum_ip_{i,t}},\qquad\text{a load is modelled if and only if }\varsigma\ge\varsigma_0=1\%.
\end{equation}
The gate acts on the measured balance share. A small balance makes the predicted voltages with and without
the extra leaf similar, so small measurement errors can change its fitted position. On the 288 development cases the statistic separates cleanly (Table~\ref{tab:gate}). The threshold is
fixed from these cases before any test, and the gate's error rates are reported by load share and kind. The spread without a
load comes from the meters' gain errors (0.3\%, 0.2\% at the PCC) and the loss estimate, so other meter classes need their
own calibration.
\begin{table}[h]
\caption{Unmetered share $\varsigma$ on the development cases (no identifier involved)}\label{tab:gate}
\centering
\begin{tabular}{@{}lrr@{}}
\toprule
Cases & Number & $\varsigma$: min / median / max\\
\midrule
No load & 72 & $-0.49\%$ / $0.03\%$ / $0.66\%$\\
Load of 2\% (three kinds) & 72 & $1.53\%$ / $2.0\%$ / $2.61\%$\\
Load of 5\% & 72 & $4.59\%$ / $5.0\%$ / $5.81\%$\\
Load of 10\% & 72 & $9.51\%$ / $10.2\%$ / $11.68\%$\\
\bottomrule
\end{tabular}
\end{table}

% ---------------------------------------------------------------------
\section{Score and labels}\label{sec:score}

\paragraph*{Score} The meter offsets $o$ are unrestricted constants, removed by $D=I_T-\Pi_{[\one]}$ (subtracting each
meter's time mean). The PCC error is random and shared by all meters in a slot: with variance $\sigma_0^2$ against
$\sigma_e^2$ for a meter, a slot's error covariance is $\Sigma=\sigma_e^2(I_n+\rho\one\one\tp)$, $\rho=\sigma_0^2/\sigma_e^2$,
taken as the ratio of the means of $V_0^4$ and $V^4$. The weighting $S_\gamma=I_n-\gamma\one\one\tp/n$ with $\gamma=1-(1+n\rho)^{-1/2}$ satisfies
$S_\gamma\tp S_\gamma=\sigma_e^2\Sigma^{-1}$ (generalised least squares), and $M=D\otimes S_\gamma$. The score is
\begin{equation}\label{eq:score}
H(F)=\min_{\theta}\frac{\|M(y-\Phi_F\theta)\|^2}{2\hat\sigma^2}+\beta\,k(F),\qquad\beta=\tfrac12p_s\ln N-\lambda,
\end{equation}
with $p_s\in\{1,2\}$ selected by an empirical check on reactive-power variation relative to active power.
This common penalty setting approximates the number of distinguishable coefficients per line.
\emph{Noise variance.} $\hat\sigma^2$ is the residual variance under $M$ of the saturated model, in which every meter sees
every leaf's flow with its own coefficient,
\begin{equation}\label{eq:sat}
y_{i,t}=\sum_{j\in\Nset}\big(R_{ij}\,p_{j,t}+X_{ij}\,q_{j,t}\big)+R_{is}\,u^P_t+X_{is}\,u^Q_t+\varepsilon_{0,t}+o_i+e_{i,t},
\qquad R=R\tp,\ X=X\tp\ \text{over the meters},
\end{equation}
with every coefficient free and the load's two terms only when a load is modelled. Every tree's prediction has this form, and
with the load's terms free $\hat\sigma^2$ does not depend on the load's place. Its degrees of freedom are the rank of $M$ minus
the numerical rank of the saturated design. It is computed once, at set-up, like the starting trees
(Appendix~\ref{app:pool}); the score, the searches and the labels hold the estimated load profile fixed. The working covariance accounts for the shared PCC error; other dependencies may affect the scores.
\begin{proposition}[Structural penalty]\label{prop:prior}
For fixed leaf count, the branching penalty $-\ln(a)\sum_v(c_v-2)$ equals a constant minus $\lambda k(F)$,
where $c_v$ is the number of children of internal node $v$ and $\lambda=\ln(1/a)$.
\end{proposition}
The default $a=0.141$ comes from branching statistics on 84 LV feeders and gives $\lambda=1.96$.
The same feeders informed the case generator. A paired comparison with $\lambda=0$ would isolate this empirical
adjustment. We use $m=4.6$ as a tolerance in score units. Its effect on coverage and errors requires measurements
against known groupings.

\paragraph*{Groupings} With a load, one grouping of meters can appear as two supports, $M$ and $M\cup\{s\}$: the same
statement about the meters, with the load inside the line or not. Write $\pi(C)=C\cap\Nset$ and call
$\Gset(F)=\{\pi(C):C\in F,\ 2\le|\pi(C)|\le n-1\}$ the \emph{groupings} of $F$. Labels, anchors and probing concern
groupings; the load's place is reported separately. Without a load, groupings are lines.

\paragraph*{Labels} For a pool $\Omega$ and $H^*=\min_FH(F)$, let
$\Sset_m=\{F\subseteq\Omega\cup\Bset\ \text{laminar}:\ H(F)\le H^*+m\}$. For a tree $T$ and a grouping $M$, the \emph{flip
problem} asks for the best tree that decides $M$ the other way,
\begin{equation}\label{eq:flip}
H^{\rm f}_M(T)=\min\{H(F):\ F\subseteq\Omega\cup\Bset\ \text{laminar},\ [M\in\Gset(F)]\neq[M\in\Gset(T)]\},
\end{equation}
which for $M\in\Gset(T)$ is the tree MIQP (\ref{eq:master}) with every support of grouping $M$ removed from the pool, and
for $M\notin\Gset(T)$ the same MIQP with a support of grouping $M$ required. Its
\emph{gap} is
\begin{equation}\label{eq:gap}
e_M(T)=H^{\rm f}_M(T)-H(T)\ \in\ \big[\,L_M-H(T),\ U_M-H(T)\,\big],
\end{equation}
with $L_M$ the solver's lower bound and $U_M$ the score of the best flipped tree found, the \emph{rival} $F_M$.
\begin{definition}[Labels]
Relative to a feasible reference $T$ and its pool, $M$ is \emph{decided} if the search excludes every flipped tree up
to $H(T)+m$. It is \emph{open} when a refitted feasible rival reaches that threshold, and \emph{unsettled} when neither
conclusion is established. The implementation uses a small numerical tolerance at the cutoff.
\end{definition}
\begin{proposition}[What a label certifies]\label{prop:label}
For a decided $M$, every tree in $\Sset_m$ agrees with $T$ on $M$. For an open $M$ with rival $F_M$, a sufficient
condition for disagreement within $\Sset_m$ is
$\max\{H(T),H(F_M)\}\le L+m$, where $L$ is a valid lower bound on the unrestricted minimum in this same pool.
\end{proposition}
In particular a proven optimal reference suffices. Otherwise an open label reports a rival relative to the reference;
the sufficient condition can be checked when a lower bound on the same domain is available. Expanding the pool or
changing the data calls for a new test. The implementation tests every grouping of the reference; for groupings
outside it, it reports successful local flips and leaves the others unclassified.

\paragraph*{Local flips} Removing the lines of grouping $M$ from $T$, or swapping in a pool line $D$ whose grouping crosses
$M$ (the lines of $T$ that cross $D$ removed), gives flipped trees without solving an MIQP:
\begin{equation}\label{eq:local}
m_M(T)=\min\Big\{H(T\setminus\pi^{-1}(M)),\ \min_{D:\ \pi(D)\cross M}H(T_{|D}\cup\{D\})\Big\}-H(T),\qquad
\ell_D(T)=H(T)-H(T_{|D}\cup\{D\}),
\end{equation}
with $T_{|D}=\{E\in T:\neg(E\cross D)\}$.
\begin{lemma}[Local flips give upper bounds only]\label{lem:local}
$e_M(T)\le m_M(T)$ for $M\in\Gset(T)$, and $e_{\pi(D)}(T)\le-\ell_D(T)$ for $\pi(D)\notin\Gset(T)$. Local flips can prove a
grouping open, never decided.
\end{lemma}

\paragraph*{Why a true line is decided} Labels and anchors are useful only if removing a true line costs a lot. Let the data
follow (\ref{eq:model}) with the true tree $T^0$ and coefficients inside the bounds. After $M$, write the data as
$\tilde y=\mu+\varepsilon$: $\mu$ is the true voltage signal and $\varepsilon$ the noise. For a tree $F$, let $P_F$ project
onto everything $F$ can predict (the span of its columns), and let $d_F=\|\mu-P_F\mu\|$: the part of the true signal that $F$
cannot reproduce with any coefficients. $\mathrm{RSS}(F)$ is the residual of $F$'s bounded fit, the numerator of
(\ref{eq:score}).
\begin{theorem}[Removing a true line costs a lot]\label{thm:cost}
(a) If $F$ lacks a line of $T^0$, then, with $\xi_F=\langle\mu-P_F\mu,\ \varepsilon\rangle/d_F$ for $d_F>0$ and $\xi_F=0$ otherwise,
\[
\mathrm{RSS}(F)-\mathrm{RSS}(T^0)\ \ge\ d_F^2-2\,d_F\,|\xi_F|-\|P_F\varepsilon\|^2 .
\]
(b) If $F$ contains every line of $T^0$ and the unbounded least-squares fit of $T^0$ lies within the bounds, then
$0\le\mathrm{RSS}(T^0)-\mathrm{RSS}(F)\le\|(P_F-P_{T^0})\varepsilon\|^2$.
\end{theorem}
The first bound measures the part of the voltage response a rival cannot reproduce. The second bounds the benefit of
adding redundant lines by the noise projected onto their additional directions. If different branch loads vary together,
$d_F$ may be small even over a long record. More slots strengthen separation when they add distinguishable signal.

\begin{corollary}[When the search finds the true tree]\label{cor:truth}
Suppose the pool contains $T^0$, the unbounded least-squares fit of $T^0$ lies within the bounds, and (A) every tree $F$ of
the pool that lacks a line of $T^0$ has
$\big(d_F^2-2d_F|\xi_F|-\|P_F\varepsilon\|^2\big)/(2\hat\sigma^2)>m+\beta\,\big(k(T^0)-k(F)\big)$, and (B) every tree
$F\supsetneq T^0$ of the pool has $\|(P_F-P_{T^0})\varepsilon\|^2<2\hat\sigma^2\beta\,\big(k(F)-k(T^0)\big)$. Then $T^0$ is the unique best tree over the pool;
when the flip tests finish, every grouping of $T^0$ is decided; and if every anchor node (including the root) has at most $K$ units, the anchored
search of Proposition~\ref{prop:anchor} still contains $T^0$, so restricting the search to it loses nothing.
\end{corollary}
These sufficient conditions concern all competing trees in the specified pool and involve the unknown true signal and
disturbance. Development checks instead measure errors and coverage against known topology.

% ---------------------------------------------------------------------
\section{Identification}\label{sec:id}

\paragraph*{Tree MIQP} For a pool $\Omega$ of supports, with a binary $b_C$ per support,
\begin{equation}\label{eq:master}
\begin{aligned}
\min_{b,r,x}\ &\frac{\|M(y-\Phi(r,x))\|^2}{2\hat\sigma^2}+\beta\sum_{C\in\Omega}b_C\\
\text{s.t.}\ \ &0\le r_C\le\bar r_Cb_C,\quad0\le x_C\le\bar x_Cb_C\quad(C\in\Omega),\qquad b_C+b_D\le1\quad(C\cross D),\\
&\textstyle\sum_{C\in\Omega,\,C\subsetneq W}b_C\le|W|-2\quad(W\in\Omega\cup\{\Nset^+\}),\qquad
\sum_{C\in\Omega,\ \pi(C)=A}b_C\ge1\quad(A\in\Aset),\\
&\textstyle\sum_{C\ni i}r_C\le\bar R_i,\quad\sum_{C\ni i}x_C\le\bar X_i\quad(i\in\Nset;\ \text{optional}),
\end{aligned}
\end{equation}
with $\Bset$ always on. The count cut holds for every tree (a tree on $|W|$ leaves has at most $|W|-2$ lines strictly inside
$W$). The fourth constraint keeps every anchor's grouping, with the load inside or not. At a score level $\ell$, the optional bounds $\bar R_i,\bar X_i$ come from a continuous relaxation of the
per-meter path sums; $\bar r_C=\min(\bar\theta,\min_{i\in\pi(C)}\bar R_i)$ and similarly for $x$.
For nonnegative line prices they preserve every tree at or below $\ell$. An incumbent in this range also preserves
the original minimum. In a threshold-only flip test, a restricted dual bound $b$ gives the conservative original-domain
bound $\min(b,\ell)$; exclusion of the threshold is recorded separately. At residual equality the implementation retains
the original bounds. Each main tree
solve first runs the local search to convergence: no removal of one line, and no swap of one pool line in (the lines it
crosses removed), lowers the score, so the value of every candidate line has been checked. The MIQP starts from that tree,
and its tree, when better, is searched to convergence too. Its claims are checked against the trees known: a bound above
a known tree's score, or an optimum that one move still improves, is a numerical failure, and such a solve gives no bound
and decides nothing. Each coefficient enters the solver in units of its column's size, so that the quadratic has a unit
diagonal, and the solver runs at its default tolerances: unscaled, with the tolerances forced to $10^{-9}$, it declared
feasible relaxations infeasible, and claimed false optima, on feeders of 47 to 60 meters.
\begin{proposition}[The load's path needs no constraints]\label{prop:path}
For a set of meters $C$ write $C^1=C\cup\{s\}$. Over $\Nset^+$ the crossing constraints of (\ref{eq:master}) are exactly:
(i) the crossings of the meter sets; (ii) one place: $b_{C^1}+b_{D^1}\le1$ unless $C\subseteq D$ or $D\subseteq C$; (iii) a
line above a line that carries the load carries it too: $b_{C^1}+b_D\le1$ for meter sets $C\subsetneq D$. So the load's path
needs neither extra constraints nor products of binaries and coefficients.
\end{proposition}

\paragraph*{Node pricing} Let $v$ be an internal node of $T$ with children $A_1,\dots,A_{c_v}$ (supports or leaves).
\begin{lemma}[Node decomposition]\label{lem:node}
A support $C\notin T$ with at least two leaves is compatible with $T$ if and only if $C=\bigcup_{j\in S}A_j$ for a unique
internal node $v$ and a set $S$ of its children with $2\le|S|<c_v$.
\end{lemma}
\begin{equation}\label{eq:price}
g_v(T)=\max_{S:\ 2\le|S|<c_v}H(T)-H\Big(T\cup\Big\{\textstyle\bigcup_{j\in S}A_j\Big\}\Big).
\end{equation}
A node with at most $K$ children is priced by fitting every union, as in the anchored universe below. A wider node is
priced by a node MIQP with a binary $z_a$
per child; the new line's coefficient enters as $u_{ab}=rz_az_b$, written exactly by the linear inequalities $u_{ab}\le r$,
$u_{ab}\le\bar\theta z_a$, $u_{ab}\le\bar\theta z_b$, $u_{ab}\ge r-\bar\theta(2-z_a-z_b)$; every coefficient of $T$ is
refitted, with the coefficients scaled to their columns as in the tree MIQP. Enumeration covers all eligible unions; a completed node MIQP can cover the same domain with one solve. A single-line
certificate requires every node to exclude a gain exceeding $10^{-7}$ in the implementation's objective
$\mathrm{crit}=2H/N$. For a wide node this requires a valid lower bound on its best augmented score. A feasible
paying line, a complete exclusion, and an unresolved node are recorded separately. The restricted tree solve has its
own domain, required anchors, incumbent, bound and gap tolerance. Two lines can improve together when neither pays
alone (Remark~\ref{rem:pairs}).
\begin{remark}\label{rem:pairs}
Four meters, fit term $\tfrac12\|R_{\rm obs}-R\|_F^2$ and price $0.6$; $R_{\rm obs}$ has diagonal $100$ and off-diagonal
$1$, except $2$ for the pairs $\{1,2\}$ and $\{3,4\}$. The star scores $4/3$, the best single line $1.4$, and the lines
$\{1,2\}$ and $\{3,4\}$ together $1.2$.
\end{remark}

\paragraph*{Anchors}
\begin{definition}[Anchors and units]
An \emph{anchor} is a decided grouping of $T$. For a laminar set of anchors $\Aset$ and each $A\in\Aset\cup\{\Nset\}$, the
\emph{units} $\Kset_A$ are the largest anchors strictly inside $A$ and the meters of $A$ outside them; $k_A=|\Kset_A|$. For
$K\ge2$ the \emph{anchored universe} is the set of meter lines
\begin{equation}\label{eq:universe}
\Omega_K(\Aset)=\Aset\ \cup\bigcup_{A:\,k_A\le K}\Big\{\textstyle\bigcup_{B\in S}B:\ S\subseteq\Kset_A,\ 2\le|S|\le k_A-1\Big\}
\ \cup\ \Omega_{>K}(\Aset),
\end{equation}
with $\Omega_{>K}(\Aset)$ the groupings of the pool compatible with $\Aset$ whose smallest enclosing member of
$\Aset\cup\{\Nset\}$ has more than $K$ units. With a load, $\Omega^+_K(\Aset)$ is every line of $\Omega_K(\Aset)$ with and
without $s$, the supports $\{i,s\}$, and $\Nset$; without a load, $\Omega^+_K=\Omega_K$.
\end{definition}
\begin{proposition}[Anchored universe]\label{prop:anchor}
Let $F$ be a tree over $\Nset^+$ that keeps every anchor ($\Aset\subseteq\Gset(F)$). Every line $C$ of $F$ has $\pi(C)$ equal
to an anchor, to a union of 2 to $k_A-1$ units of one $A\in\Aset\cup\{\Nset\}$, to a single meter, or to $\Nset$. Hence, if
every $k_A\le K$, (\ref{eq:master}) over $\Omega^+_K(\Aset)$ with the anchors required, solved to optimality, returns the
best tree that keeps every anchor, among all supports and every place of the load; it has $\sum_A(2^{k_A}-k_A-2)$ union
binaries (twice that with a load) instead of $2^n-n-2$.
\end{proposition}
The result is conditional: exact over the trees that keep the anchors, at nodes with at most $K$ units, for a solve that
finishes. The anchors themselves are re-tested, and may be flipped, in every round on the current data and pool. A node with
more than $K$ units keeps the pool's lines inside it and is priced by (\ref{eq:price}).

\paragraph*{Pricing, re-solve and stopping} Every line added by pricing stays in the pool used by subsequent
flip tests, even if a restricted re-solve drops it from the tree. Bounds from a restricted solve refer to its
recorded lines and required anchors.
\begin{lemma}[Re-solve and stopping]\label{lem:stop}
The best feasible score cannot rise in a repair. With fixed anchors, a closed restricted gap excludes improvements
beyond its tolerance in that solve's domain. Complete pricing on the retained tree excludes a qualifying single-line addition.
\end{lemma}
The loop reports convergence when the anchors are unchanged, pricing added no line, every node passed the single-line
check, and the restricted gap closed to tolerance. The round limit is reported separately.

\paragraph*{Load evidence} Whether a load is modelled is decided by the gate (\ref{eq:gate}). The fit reports, as a
diagnostic, how strongly the voltages show the load: with $H_0$ the score of the same data and $\hat\sigma^2$ with
$u\equiv0$, minimised over meter topologies, and $H$ minimised with the load,
\begin{equation}\label{eq:evidence}
D^*=\min H_0-\min H,\qquad L_0-U_1\ \le\ D^*\ \le\ U_0-L_1,
\end{equation}
where $[L_0,U_0]$ and $[L_1,U_1]$ bracket the two minima (solver bounds and incumbents): $\min H_0\ge L_0$ and $\min H\le U_1$
give the lower end, and likewise the upper. Both brackets come from the tree MIQP over the final pool (over its meter
groupings for $H_0$), started from $\hat T$; $U_1$ is the better of that solve and $H(\hat T)$. Evaluating $H_0$ at the load tree's own meter topology instead could only
overstate $D^*$, since that topology is one of those minimised over. The places of the load are the admissible places in
$\hat T$'s meter topology whose score is within $m$ of the best.

\begin{algorithm}[t]
\caption{\textsc{Identify}: gate, tree, anchors and labels}\label{alg:identify}
\begin{algorithmic}[1]
\Require data; $\varsigma_0$, $\lambda$, $m$, $K$, $R$
\Ensure the gate's decision; tree $\hat T$ (groupings, and with a load its path); certificate and scope; anchors; pool;
labels with gap intervals and rivals; with a load, its places and evidence interval
\State \textbf{if} the PCC power is metered: $\varsigma$ by (\ref{eq:gate}); a load is modelled iff $\varsigma\ge\varsigma_0$ \Comment{gate}
\State $M$, $\hat\sigma^2$, $\beta$; $\hat R$ by (\ref{eq:rhat}) \Comment{set-up}
\State $\Omega\gets\Omega_0$ and nine starting trees \Comment{Appendix~\ref{app:pool}}
\State $T\gets$ (\ref{eq:master}) over $\Omega$, from the converged local-search tree
\For{$r=0,1,\dots,R$}
  \State label the groupings of $T$: local flips (\ref{eq:local}), then flip tests (\ref{eq:flip}); $\Aset\gets$ the decided ones
  \State \textbf{if} the stopping rule of Lemma~\ref{lem:stop} holds, or $r=R$: \textbf{break}
  \State $T\gets$ (\ref{eq:master}) over $\Omega^+_K(\Aset)\cup\Lin(T)$, anchors required; $\Omega\gets\Omega\cup\Omega^+_K(\Aset)$
  \State pricing (\ref{eq:price}), by enumeration or the node MIQP: add paying lines to $T$ and $\Omega$; if any, (\ref{eq:master}) again; keep the better tree
\EndFor
\State $\hat T\gets T$; groupings outside $\hat T$ whose swap (\ref{eq:local}) is within $m$ are open
\State with a load: its path, its places within $m$, its evidence (\ref{eq:evidence})
\end{algorithmic}
\end{algorithm}

\paragraph*{Explanation} Line 1 decides whether there is a load; below the gate the rest runs on the meters alone. Lines 2--4
build a first tree from cheap candidate lines. Lines 5--9 are the core: decided groupings become anchors; the anchored
universe offers every line compatible with them at nodes that are small enough (Proposition~\ref{prop:anchor}); pricing
searches the wider nodes and checks the result against any single line, enumerating nodes with at most $K$ children and
solving the node MIQP at wider ones; the loop stops by Lemma~\ref{lem:stop} or after $R$
searches, and the output says which. Every round labels all groupings of the current tree on the current pool, anchors
included, so the final labels are those of the final tree and pool. Lines 10--11 report what remains open and, with a load,
where it sits and how strongly the voltages show it.

% ---------------------------------------------------------------------
\section{Probing}\label{sec:probe}

The operator can buy signed active-power blocks at chosen meters from an aggregator of household batteries and EV chargers,
each held for one market hour and sealed one hour ahead. Probing works on the pool, the leader and the open alternatives; it
never lists whole trees. A block that changes the flows on one side of a branching point makes the leader's and the rival's
predictions differ more: this is what probing buys.

\paragraph*{Values used in probing} For an open alternative $M$ with rival $F_M$, at the gate of hour $h$:
\begin{itemize}
\item the gap interval (\ref{eq:gap}) and the \emph{deficit} $\underline\delta_M=m-\bar e_M$ with $\bar e_M=U_M-H(\hat L)\ge
e_M$: a lower estimate of what the alternative still needs; unsettled groupings are not targets (probing cannot repair a
solver's time limit) and are labelled again at each gate;
\item the fitted parameters of the leader and the rival and their covariances $P=\hat\sigma^2(\Phi\tp\Phi)^{-1}$ on the
coordinates not at a bound;
\item the predicted voltage drops of hour $h+1$ under each tree after $M$, so that $\hat\sigma^2I$ is their noise
covariance, from the day-ahead forecast of the meters' load (the mean of the same hour over the data) plus the block $u$;
their difference is $d_M(u)=d_M(0)+G_Mu$, and
\begin{equation}\label{eq:sep}
q_M(u)=d_M(u)\tp\Sigma_M^{-1}d_M(u),\qquad \Sigma_M=\hat\sigma^2I+\tfrac12\Phi_{\hat L}P_{\hat L}\Phi_{\hat L}\tp+\tfrac12\Phi_{F_M}P_{F_M}\Phi_{F_M}\tp,
\end{equation}
with $\Sigma_M$ held fixed at the forecast;
\item the feasible blocks $\Uset_{h+1}$ (battery and EV headroom, the voltage band for the leader \emph{and every rival}, the
transformer rating, the price caps) and the price $c(u)=c^\uparrow\sum_i\max(u_i,0)+c^\downarrow\sum_i\max(-u_i,0)$. A block
acts at one meter with one sign, at a quarter, a half, three quarters or all of the largest size these allow; the band keeps
$z\hat\sigma/2$ of margin (one reading of $V$).
\end{itemize}
\begin{lemma}[Expected predictive log-likelihood ratio]\label{lem:evidence}
With the two fitted predictive means $\mu_{\hat L}(u)$, $\mu_{F_M}(u)$ and the covariance $\Sigma_M$ held fixed, let
\begin{equation}\label{eq:llr}
\ell_M(Y;u)=\ln\frac{\mathcal N\big(Y;\mu_{\hat L}(u),\Sigma_M\big)}{\mathcal N\big(Y;\mu_{F_M}(u),\Sigma_M\big)} .
\end{equation}
If the next hour's data are $\mathcal N(\mu_{\hat L},\Sigma_M)$, then $\ell_M\sim\mathcal N(\tfrac12q_M,q_M)$.
\end{lemma}
$\tfrac12q_M$ is a forecast of the evidence, not of the refitted score's change: the score refits both trees, its price grows
with $N$, and refitting can move it either way.
\begin{lemma}[Reach]\label{lem:reach}
If $G_Mv=0$ for every feasible block direction $v$, no feasible block changes $q_M$.
\end{lemma}
$G_M$ must be computed from the two fitted trees, not from their structure alone. For example, two trees that differ only in
whether the line $\{1,2\}$ carries the load can fit the history equally when the load's profile equals $p_1+p_2$ (the first
with $r=2$ and no load, the second with $r=1$ and the load); a block $a$ at meter 1 then gives $2a$ against $a$, so a block
does separate them.

\paragraph*{Target set and objective} At the start of the day, $\Tset$ is the set of open groupings, and the load's place
when the second-best admissible place is within $m$ of the best (that place is its rival), whose $G$ moves on some feasible
block direction; it stays fixed during the day. At every gate each target is relabelled on the current data and pool; a
target outside the tree that no local swap shows open gets the flip test (\ref{eq:flip}) with its grouping required, since a
local swap cannot decide it (Lemma~\ref{lem:local}); its best swap is the known tree that the test's claim is checked
against. With
$\kappa(\delta,q)=\mathbb E\min(\delta,Z)$ for $Z\sim\mathcal N(q/2,q)$, that is
$\kappa=\tfrac q2-(\tfrac q2-\delta)\,\mathrm F_0\big(\tfrac{q/2-\delta}{\sqrt q}\big)-\sqrt q\,\mathrm f_0\big(\tfrac{q/2-\delta}{\sqrt q}\big)$
with $\mathrm F_0$ and $\mathrm f_0$ the standard normal distribution and density ($\kappa=\min(\delta,0)$ at $q=0$), the
block for hour $h+1$ is
\begin{equation}\label{eq:block}
u^\ast=\arg\max_{u\in\Uset_{h+1}}\ J_h(u),\qquad
J_h(u)=\psi\sum_{M\in\Tset\ \text{open}}\Big[\kappa\big(\underline\delta_M,q_M(u)\big)-\kappa\big(\underline\delta_M,q_M(0)\big)\Big]-c(u),
\end{equation}
bought only if $J_h(u^\ast)>0$: the expected capped log-likelihood ratio (\ref{eq:llr}) gained over passive operation,
counted up to what each target still needs, against the price. It is a forecast that assumes the leader generates the data,
with both fits and the covariance held fixed. (Capping the mean, $\min(\delta,q/2)$, would overstate it: at $q=2$, $\delta=1$
it gives $1$ against $0.436$.) \emph{Choice of $\psi$:} on the development cases, Algorithm~\ref{alg:probe} runs for every
$\psi$ on a grid; with mean cost and coverage scaled to $[0,1]$, the default is the grid point farthest above the straight
line from the first point to the last (the knee). Grid and default are fixed before any probing test. \emph{Greedy, in plain
words:} each hour the rule buys the one block worth most for the next hour alone. \emph{Stopping and trust gate:} the day
ends when every target holds a fresh decided label, on that gate's data and pool, at two consecutive gates; blocks are
bought only when the leader fits as well as the meter noise allows, $s^2_{\hat L}/\sigma^2_y\le\kappa_A$, with $s^2_{\hat L}$ the
leader's residual variance after $M$ and $\sigma^2_y=4\sigma_v^2\,1.02^4$ for the meter class $\sigma_v=0.05\%$. A last gate after
hour 23 relabels on the whole day, so that the readings of the last block are used.

\begin{algorithm}[t]
\caption{\textsc{Probe}: one day of cost-aware probing (proposed)}\label{alg:probe}
\begin{algorithmic}[1]
\Require the output of Algorithm~\ref{alg:identify}; history; contract; $\psi$, $m$, $\kappa_A$
\Ensure the leader, with every grouping decided, open or unsettled
\State $\Tset\gets$ open alternatives whose $G$ moves on a feasible block direction \Comment{fixed for the day; Lemma~\ref{lem:reach}}
\For{$h=0,1,\dots,23$} \Comment{gate at the start of hour $h$; block for $h+1$ sealed now}
  \State $\hat L\gets$ Algorithm~\ref{alg:identify}, lines 5--9, on all data so far
  \State relabel every target on the current data and pool (an outside target by its flip test): intervals, deficits, rivals
  \State \textbf{if} every target has a fresh decided label at this gate and the previous one: \Return $\hat L$ with its labels
  \State $u_{h+1}\gets u^\ast$ of (\ref{eq:block}) if $s^2_{\hat L}/\sigma^2_y\le\kappa_A$ and $J_h(u^\ast)>0$, else $0$
  \State hour $h$ runs with $u_h$; add its readings to the data
\EndFor
\State lines 3--4 once more on the whole day (the last gate); \Return $\hat L$ with its labels \Comment{some may stay open}
\end{algorithmic}
\end{algorithm}

% ---------------------------------------------------------------------
\section{How the outputs are valued}\label{sec:value}

\paragraph*{What the output says} Three parts, kept apart. \emph{The load:} the gate's statistic and decision; with a load,
its path, places within $m$ and evidence interval. \emph{The structure:} the groupings of $\hat T$, the anchors, and for every
tested grouping its label, gap interval and any available rival. \emph{The computation:} the certificate's scope (single-line pricing complete, a paying line left at the repair limit, or
incomplete; and whether the restricted scope holds), the nodes priced each way, the
share of unsettled labels, the solver status of every number, and why the loop stopped.

\paragraph*{How it is scored against the truth} The gate: loads missed and no-load feeders given a load, by share and kind.
The structure: precision, recall and F1 over groupings; recall near the root; exact trees; false anchors; label coverage,
false decisions, the share of open groupings that are true, and the share unsettled under each time limit. The load: whether
its true place is within $m$ of the best. Probing: the share of targets decided by the end of the day, false decisions
against passive operation, cost, energy and time; over $\psi$, as a cost--coverage curve. Coverage and false decisions are
always reported together.

% ---------------------------------------------------------------------
\section{Contributions, limits and tests}\label{sec:tests}

\paragraph*{Contributions} The identification target is agreement on downstream meter groups when both the feeder
and one balance-derived load location may vary. A group flip removes all supports with that meter projection
(Proposition~\ref{prop:path}). Decided groups define anchor units; when every unit count is at most $K$, unions
cover all anchor-preserving trees and admissible load paths (Proposition~\ref{prop:anchor}). For larger nodes,
the implementation reports the smaller scope it searched. Theorem~\ref{thm:cost} explains which part of the
physical response separates a rival. Probing uses the remaining alternatives to rank feasible actions; its
experimental assessment is separate.

\paragraph*{Relation to earlier work} Bremer measured the additional tree length needed to lose a branch in
near-most-parsimonious trees~\cite{bremer}. Marx, Calmon and Ustun minimised classification error with one prediction
flipped to study disagreement between near-optimal models~\cite{marx}. Pengwah et al.\ studied topology recovery with
partial smart-meter coverage through impedance estimation followed by meter placement~\cite{pengwah}. Here the score uses measured
voltage and power, the tested object is a meter projection that can span several load placements, and the search
reports its coverage of alternative feeder structures.

\paragraph*{Limits} Labels depend on the gate, fit, penalty, margin, data and candidate pool. Multiple unmetered
loads and errors in measured regressors can change the fitted groups. The gate was calibrated for one meter class.
The signal-separation conditions involve the unknown true response. Probing uses fixed fits and a conditional
forecast model. The proposed assessments report coverage, errors and unresolved cases together.

\paragraph*{Relation to the registered test} The registered sealed test evaluates the predecessor (final.py at 71684ff). It
has no gate, removes the load during the search by projecting out its profiles ($D_1$, Appendix~\ref{app:pool}) and places it at the end, takes blocks from the top-level
lines of the first tree without flip tests, searches the unions of at most 9 blocks once, prices nodes with an MIQP over the
meters only, has no labels, and computes the evidence at the load tree's meter topology. Its result does not carry over to
Algorithm~\ref{alg:identify}.

\paragraph*{Development checks before a registered test} (V1) Do meter-level labels separate grouping uncertainty from the
load's place? (V2) The gate's error rates; the evidence at the load tree's topology against the two-model evidence. (V3) Scale
and time per case. (V4) How often nodes are wide or unproved, and the effect on F1. (V5) With and without the per-meter
bounds. (V6) Decided, open and unsettled shares under several time limits. (V7) Whether better gap intervals give better
probing. (V8) How grouping coverage, disagreement with truth and the unsettled share vary with the margin and candidate pool.

\paragraph*{Small exhaustive audit} Four paired inputs use four meters, 96 hourly observations, two excitation
levels and either zero or one extra leaf. The complete domains contain 26 and 236 trees, respectively. Using the same
reference and score threshold within each input, the two strongly excited cases retained all four initially decided
groupings across them. In the weakly excited cases the full domains admitted competitive trees for four and five
meter groups respectively that had no support in the initial pools. The full-domain optimum in each weak case was
the star. The strong input with the extra leaf illustrates the distinction between an identified group and a load location.
Five of the 236 trees lay within 4.6 of the full-domain optimum. All five had exactly the meter groups
$\{0,1\}$ and $\{2,3\}$ (zero-based indices), with four distinct sets of internal lines carrying the extra leaf.
Two such trees place it below group $\{2,3\}$ and outside both groups, with respective score increments 0
and 2.83. These cases check implementation and candidate coverage at this scale.

% ---------------------------------------------------------------------
\appendices
\section{Proofs}

The signal-separation result is deterministic conditional on the linearised model. The other proofs specify search
domains, label semantics and the forecast used for probing.

\begin{IEEEproof}[Proof of Theorem~\ref{thm:cost}]
Two facts about least squares. (F1) A tree that contains every line of $T^0$ can use the true coefficients (and zero on its
other lines), which leave the residual $\tilde y-\mu=\varepsilon$; its best fit is no worse, so its RSS is at most
$\|\varepsilon\|^2$. (F2) No fit of $F$, bounded or not, beats the unbounded least-squares fit, whose residual is
$Q_F\tilde y$ with $Q_F=I-P_F$; so $\mathrm{RSS}(F)\ge\|Q_F\tilde y\|^2$.
(a) By (F1) for $T^0$ and (F2) for $F$, $\mathrm{RSS}(F)-\mathrm{RSS}(T^0)\ge\|Q_F(\mu+\varepsilon)\|^2-\|\varepsilon\|^2$.
Expanding the square, and using $\|Q_F\varepsilon\|^2=\|\varepsilon\|^2-\|P_F\varepsilon\|^2$ (Pythagoras),
\[
\|Q_F(\mu+\varepsilon)\|^2-\|\varepsilon\|^2=\|Q_F\mu\|^2+2\langle Q_F\mu,\varepsilon\rangle-\|P_F\varepsilon\|^2
=d_F^2+2d_F\xi_F-\|P_F\varepsilon\|^2\ \ge\ d_F^2-2d_F|\xi_F|-\|P_F\varepsilon\|^2 ,
\]
since $Q_F$ is symmetric and $Q_F^2=Q_F$, so that $\langle Q_F\mu,Q_F\varepsilon\rangle=\langle Q_F\mu,\varepsilon\rangle$.
(b) $F$ can reproduce any fit of $T^0$ by setting its extra lines' coefficients to zero, so the difference is at least $0$.
The unbounded fit of $T^0$ lies within the bounds, so it is also its bounded fit; $\mu$ lies in the span of $T^0$'s columns, so
$\mathrm{RSS}(T^0)=\|Q_{T^0}\tilde y\|^2=\|Q_{T^0}\varepsilon\|^2$. $\mu$ lies in the span of $F$'s columns too, so (F2) gives
$\mathrm{RSS}(F)\ge\|Q_F\varepsilon\|^2$. Subtracting, and using Pythagoras twice,
$\|Q_{T^0}\varepsilon\|^2-\|Q_F\varepsilon\|^2=\|P_F\varepsilon\|^2-\|P_{T^0}\varepsilon\|^2=\|(P_F-P_{T^0})\varepsilon\|^2$:
the span of $T^0$'s columns lies in that of $F$'s, so $P_F-P_{T^0}$ projects onto the part of $F$'s span orthogonal to
$T^0$'s.
\end{IEEEproof}

\begin{IEEEproof}[Proof of Corollary~\ref{cor:truth}]
The score is $H=\mathrm{RSS}/(2\hat\sigma^2)+\beta k$. Take a tree $F\ne T^0$ of the pool. If $F$ lacks a line of $T^0$,
Theorem~\ref{thm:cost}(a) and (A) give $H(F)-H(T^0)>m$. Otherwise $F$ contains every line of $T^0$ and $k(F)-k(T^0)\ge1$ more, and Theorem~\ref{thm:cost}(b) and (B) give
$H(T^0)-H(F)\le\|(P_F-P_{T^0})\varepsilon\|^2/(2\hat\sigma^2)-\beta\,\big(k(F)-k(T^0)\big)<0$. So $T^0$ is the unique best
tree. Every tree that flips a grouping of $T^0$ lacks the line of that grouping, so it scores more than $H(T^0)+m$, and a
finished flip test proves the grouping decided. The anchors are groupings of $T^0$, so $T^0$ keeps them, and by
Proposition~\ref{prop:anchor} it lies in the anchored universe.
\end{IEEEproof}

\begin{IEEEproof}[Proof of Proposition~\ref{prop:prior}]
A tree with $n^+$ leaves (the load counts as one) and $k+1$ internal nodes has $n^++k$ edges, one from each node except the
root to its parent, so $\sum_vc_v=n^++k$ and $\sum_v(c_v-2)=n^++k-2(k+1)=n^+-k-2$. Hence
$-\ln(a)\sum_v(c_v-2)=\lambda(n^+-2)-\lambda k$.
\end{IEEEproof}

\begin{IEEEproof}[Proof of Proposition~\ref{prop:label}]
Exclusion implies $H(F)>H(T)+m\ge H^*+m$ for every flipped tree. For the second statement,
$L\le H^*$ puts both feasible trees in $\Sset_m$; their groupings disagree by construction.
\end{IEEEproof}

\begin{IEEEproof}[Proof of Lemma~\ref{lem:local}]
Each tree in (\ref{eq:local}) is laminar and decides the grouping the other way (a swapped-in $D$ whose meters cross $M$
removes every line of grouping $M$, since those lines cross $D$), so it is a feasible point of (\ref{eq:flip}), and the
minimum is at most its score. A feasible point bounds a minimum from above, never from below.
\end{IEEEproof}

\begin{IEEEproof}[Proof of Proposition~\ref{prop:path}]
$C^1$ and $D^1$ both contain $s$, so they are never disjoint: they are compatible exactly when one contains the other, that
is when $C\subseteq D$ or $D\subseteq C$. For $C^1$ and $D$ (without $s$): $C^1\subseteq D$ is impossible, $D\subseteq C^1$
holds exactly when $D\subseteq C$, and they are disjoint exactly when $C\cap D=\emptyset$; so they cross exactly when
$C\subsetneq D$ or $C$ and $D$ cross. Supports without $s$ cross as meter sets.
\end{IEEEproof}

\begin{IEEEproof}[Proof of Lemma~\ref{lem:node}]
If $C$ is a union of 2 to $c_v-1$ children of $v$, it is nested with or disjoint from every support of $T$, and equals none.
Conversely, let $v$ be the smallest support of $T$ containing $C$. No child of $v$ contains $C$, so each child is inside $C$
or disjoint from it; the children partition $v$, so $C$ is a union of children, at least two and not all.
\end{IEEEproof}

\begin{IEEEproof}[Proof of Proposition~\ref{prop:anchor}]
Removing $s$ keeps inclusions and disjointness, so $\pi(C)$ is compatible with every anchor. If $2\le|\pi(C)|\le n-1$, let
$A$ be the smallest member of $\Aset\cup\{\Nset\}$ containing $\pi(C)$. Each unit of $A$ is inside $\pi(C)$ or disjoint from
it: a meter trivially, and an anchor unit $B$ because $\pi(C)\subsetneq B$ would contradict the choice of $A$. So $\pi(C)$ is
a union of units of $A$: one unit makes it an anchor, all of them make it $A$. Otherwise $\pi(C)$ is a single meter ($C$ is
$\{i,s\}$) or $\Nset$. In every case $C$ is $\pi(C)$ or $\pi(C)\cup\{s\}$, which is in $\Omega^+_K(\Aset)$ when every
$k_A\le K$. A node with $k$ units has $2^k-k-2$ unions of 2 to $k-1$ units.
\end{IEEEproof}

\begin{IEEEproof}[Proof of Lemma~\ref{lem:stop}]
A newly added line is retained only when it lowers the score, and the re-solve keeps the better feasible tree.
The feasible score minus a valid lower bound bounds the improvement left in a restricted solve. At each node,
enumeration or a valid bound on the best gain excludes a qualifying single addition to the retained tree.
\end{IEEEproof}

\begin{IEEEproof}[Proof of Lemma~\ref{lem:evidence}]
Both densities have covariance $\Sigma_M$, so the terms quadratic in $Y$ cancel and, with $d=\mu_{\hat L}-\mu_{F_M}$,
$\ell_M=d\tp\Sigma_M^{-1}\big(Y-\tfrac12(\mu_{\hat L}+\mu_{F_M})\big)$. If $Y\sim\mathcal N(\mu_{\hat L},\Sigma_M)$, the
bracket has mean $\tfrac12d$ and covariance $\Sigma_M$, so $\ell_M$ is Gaussian with mean $\tfrac12d\tp\Sigma_M^{-1}d=
\tfrac12q_M$ and variance $d\tp\Sigma_M^{-1}\Sigma_M\Sigma_M^{-1}d=q_M$.
\end{IEEEproof}

\begin{IEEEproof}[Proof of Lemma~\ref{lem:reach}]
$d_M(u)=d_M(0)+G_Mu$ and $\Sigma_M$ is fixed, so $q_M(u)$ depends on $u$ only through $G_Mu$, which is zero for every
feasible block.
\end{IEEEproof}

\section{Initial candidate lines used in this work}\label{app:pool}

Any tree builders can supply $\Omega_0$; this work uses nine. (1) A sensitivity matrix over the meters without tree structure, fitted to the
active-power part of (\ref{eq:model}) by a convex QP solved with Gurobi. With a load, its terms $R_{is}u^P_t+X_{is}u^Q_t$ are
free for each meter, as in (\ref{eq:sat}); minimising over them first replaces $D$ by $D_1=I_T-\Pi_{[\one,\,u^P,\,u^Q]}$
(least squares in two steps), and $D_1=D$ without a load. $D_1$ is used only here and in $\hat\sigma^2$:
\begin{equation}\label{eq:rhat}
\hat R=\arg\min_R\ \big\|(D_1\otimes S_\gamma)\,\mathrm{vec}(Y-PR)\big\|^2\quad\text{s.t.}\quad R=R\tp\ge0,\quad
R_{ii}-R_{ij}\ge\epsilon\ \ (i\ne j),
\end{equation}
with $\epsilon$ a strict order margin, $10^{-6}$ of the median diagonal of the unconstrained fit.
(2) Nine hierarchies over the meters: recursive neighbour joining, neighbour joining (with the PCC as an observed leaf) and
recursive grouping (equalities to $10^{-9}$ of the median depth) on the shared paths $\hat R/\bar d$, $\bar d$ the mean
positive path distance;
UPGMA on the path distance $d_{ij}=\hat R_{ii}+\hat R_{jj}-2\hat R_{ij}$ and Ward linkage on $\sqrt d$; average-linkage
clusterings on the correlation distance between meters of the rows of $\hat R$, the drops $y$, the drops
$(D_1\otimes S_\gamma)y$, and the residuals after fitting the service cables and the trunk alone. (3) With a load, it is placed in each of
the nine trees at its best admissible place; these trees start the first solve. (4) $\Omega_0$ is every cluster with 2 to
$n-1$ meters in any hierarchy, and with a load each such cluster with and without $s$, the supports $\{i,s\}$ and $\Nset$,
so $|\Omega_0|\le18(n-2)+n+1$.

\section{Settings}

\begin{table}[h]
\caption{Settings of Algorithm~\ref{alg:identify} (development; values from the registered predecessor where it has them)
and of probing}\label{tab:settings}
\centering
\begin{tabular}{@{}p{0.36\textwidth}p{0.56\textwidth}@{}}
\toprule
Setting & Value\\
\midrule
Gate $\varsigma_0$ & 1\% of the metered mean power (prior setting; Table~\ref{tab:gate})\\
$\lambda$; $m$; $\bar\theta$ & 1.96; 4.6 (half-BIC units); 2 pu\\
$K$; anchored searches $R$ & 9 units (children in pricing); at most 2, then the stopping rule of Lemma~\ref{lem:stop} or the limit\\
Tree MIQP & 120 s up to 35 meters, 240 s beyond, always started from the converged local-search tree; coefficients
scaled to their columns; Gurobi's default tolerances; per-meter bounds on, at a quarter of the limit at most (off in the ablation)\\
Flip test & cutoff $H(T)+m$, local flips first; limit at least that of the tree MIQP; the share at the limit is reported\\
Local search & a first pass of 20 s per start (40 s beyond 35 meters) from the empty tree and the builder trees; the
best start, and the MIQP's tree when better, to convergence\\
Pricing & enumeration up to $K$ children; node MIQP beyond, 60 s; at most 10 additions per round\\
Load evidence & the no-load model over the meter groupings of the pool and the load model over the pool, same limits\\
Solver & Gurobi, one thread per case\\
\midrule
Probing & $c^\downarrow=4$, $c^\uparrow=1.2$ CNY/kWh; caps 20 CNY/h, 100 CNY/day; band $[0.96,1.035]$, $z=2.33$;
$\kappa_A=1.2$, $\sigma_v=0.05\%$; blocks at one meter, $\frac14$ to all of the largest feasible size; forecast: the mean of the
same hour; $\psi$ at the knee of the development cost--coverage curve\\
\bottomrule
\end{tabular}
\end{table}

\begin{thebibliography}{3}
\bibitem{bremer} K.~Bremer, Branch support and tree stability, \emph{Cladistics}, vol.~10, pp.~295--304, 1994.
\href{https://doi.org/10.1111/j.1096-0031.1994.tb00179.x}{Publisher record}.
\bibitem{marx} C.~Marx, F.~Calmon and B.~Ustun, Predictive multiplicity in classification, \emph{Proc. ICML},
pp.~6765--6774, 2020. \href{https://proceedings.mlr.press/v119/marx20a.html}{Proceedings}.
\bibitem{pengwah} A.~B.~Pengwah et al., Topology identification of distribution networks with partial smart meter
coverage, \emph{IEEE Trans. Power Delivery}, 2024. \href{https://doi.org/10.1109/TPWRD.2024.3354292}{Publisher record}.
\end{thebibliography}
\end{document}
